\documentclass[conference]{IEEEtran}
\IEEEoverridecommandlockouts
% The preceding line is only needed to identify funding in the first footnote. If that is unneeded, please comment it out.
\usepackage{cite}
\usepackage{amsmath,amssymb,amsfonts}
\usepackage{algorithmic}
\usepackage{graphicx}
\usepackage{textcomp}
\usepackage{xcolor}
\usepackage{booktabs}
\usepackage{url}
\def\BibTeX{{\rm B\kern-.05em{\sc i\kern-.025em b}\kern-.08em
    T\kern-.1667em\lower.7ex\hbox{E}\kern-.125emX}}
\begin{document}

\title{Eliminating Nested Scheduling in ROS 2: OS-Native Callback Dispatch on RTEMS Real-Time Systems
% \thanks{Identify applicable funding agency here. If none, delete this.}
}

% \author{\IEEEauthorblockN{1\textsuperscript{st} Given Name Surname}
% \IEEEauthorblockA{\textit{dept. name of organization (of Aff.)} \\
% \textit{name of organization (of Aff.)}\\
% City, Country \\
% email address or ORCID}
% \and
% \IEEEauthorblockN{2\textsuperscript{nd} Given Name Surname}
% \IEEEauthorblockA{\textit{dept. name of organization (of Aff.)} \\
% \textit{name of organization (of Aff.)}\\
% City, Country \\
% email address or ORCID}
% \and
% \IEEEauthorblockN{3\textsuperscript{rd} Given Name Surname}
% \IEEEauthorblockA{\textit{dept. name of organization (of Aff.)} \\
% \textit{name of organization (of Aff.)}\\
% City, Country \\
% email address or ORCID}
% \and
% \IEEEauthorblockN{4\textsuperscript{th} Given Name Surname}
% \IEEEauthorblockA{\textit{dept. name of organization (of Aff.)} \\
% \textit{name of organization (of Aff.)}\\
% City, Country \\
% email address or ORCID}
% \and
% \IEEEauthorblockN{5\textsuperscript{th} Given Name Surname}
% \IEEEauthorblockA{\textit{dept. name of organization (of Aff.)} \\
% \textit{name of organization (of Aff.)}\\
% City, Country \\
% email address or ORCID}
% \and
% \IEEEauthorblockN{6\textsuperscript{th} Given Name Surname}
% \IEEEauthorblockA{\textit{dept. name of organization (of Aff.)} \\
% \textit{name of organization (of Aff.)}\\
% City, Country \\
% email address or ORCID}
% }

\maketitle

\begin{abstract}
ROS 2's Executor introduces an independent scheduling layer at the middleware layer, forming a nested scheduling structure with the operating system's thread scheduling. 
This design leads to priority inversion, nondeterministic callback dispatch order, and increased complexity in end-to-end response time analysis, as pointed out in several RTSS/ECRTS/RTAS papers. 
Existing work (such as CallbackIsolatedExecutor and EventsExecutor) attempts to alleviate this problem on Linux, but is limited by the non-real-time nature of the Linux kernel. This paper proposes a fundamental solution: porting ROS 2 to the RTEMS 6.1 real-time operating system and completely returning callback scheduling control to the operating system kernel through RTExecutor. 
RTExecutor consists of TimeManager (hardware tick interrupt-driven timer callbacks) and EventManager (RTEMS task scheduling subscription/service callbacks), allowing each ROS 2 callback to be directly mapped to an RTEMS real-time task with an independent priority, uniformly managed by the RTEMS kernel's fixed-priority preemptive scheduler. 
We completed the full porting of ROS 2 to RTEMS on the ARM realview\_pbx\_a9\_qemu platform. 
Experiments show that: (1) the timer callback jitter is $\leqslant 2$ ticks (2ms), which is significantly better than the software timer on Linux; (2) after eliminating the middleware layer scheduling, the callback response time can be directly analyzed by the classic real-time scheduling theory (rate-monotonic / fixed-priority preemptive scheduling) without considering the nested scheduling model; (3) intra-process zero-copy communication works in conjunction with RTEMS task scheduling to achieve end-to-end deterministic data flow.
\end{abstract}

\begin{IEEEkeywords}
ROS 2, RTEMS, real-time scheduling, nested scheduling elimination, callback dispatch, priority inversion, hardware timer interrupts
\end{IEEEkeywords}

\section{Introduction}

Component-based middleware built on the publish-subscribe paradigm has become the de facto standard for developing cyber-physical systems (CPS), including autonomous vehicles, industrial robots, and surgical systems.
The Robot Operating System~2 (ROS~2) is currently the most widely adopted such middleware, providing a rich ecosystem of libraries and tools for perception, planning, and control.

A key requirement for these systems is \emph{real-time determinism}: sensor data must be processed within bounded latency, control commands must be issued at precise intervals, and high-priority callbacks must preempt lower-priority ones.
Unfortunately, the standard ROS~2 execution model falls short of these requirements.
The \texttt{rclcpp} Executor implements its own callback scheduling---a FIFO or round-robin dispatch of ready callbacks from an internal wait-set---on top of the operating system's thread scheduler.
This creates a \emph{nested scheduling} structure~\cite{casini2019rtas}, in which two independent schedulers operate at different abstraction levels without coordinating priority information.

The consequences are well documented: \emph{unbounded priority inversion} occurs when a low-priority callback occupying the Executor thread blocks a high-priority one~\cite{casini2019rtas}; \emph{non-deterministic dispatch order} arises from the wait-set mechanism~\cite{sun2024rtss}; \emph{timing anomalies} appear where adding resources actually increases latency~\cite{voronov2024rtas}; and \emph{response-time analysis} becomes intractable because both scheduling layers must be modeled simultaneously~\cite{tang2020rtss,blass2021rtss}.

Recent work has attempted to mitigate these problems within the Linux ecosystem.
CallbackIsolatedExecutor~\cite{ishikawa2025rtas} maps each callback to a dedicated \texttt{SCHED\_FIFO} thread.
EventsExecutor~\cite{teper2024rtss} bridges ROS~2 with classic fixed-priority scheduling.
PiCAS~\cite{choi2021rtas} implements priority-driven chain-aware scheduling.
However, all these approaches are fundamentally limited by the Linux kernel: CFS does not provide deterministic guarantees, interrupt latency remains non-deterministic even with PREEMPT\_RT, and the overhead of user--kernel mode switches is significant.

We propose a fundamentally different approach: \emph{eliminate the middleware scheduling layer entirely} by porting ROS~2 to a true real-time operating system and returning all scheduling decisions to the OS kernel.
Specifically, we make the following contributions:

\begin{enumerate}
    \item \textbf{Complete porting of ROS~2 to RTEMS~6.1 on ARM.} We ported the full \texttt{rclcpp} stack---including the intra-process zero-copy communication path---to the RTEMS real-time operating system running on the ARM \texttt{realview\_pbx\_a9\_qemu} platform. This is the first full ROS~2 port to an open-source RTOS that supports the complete \texttt{rclcpp} API.

    \item \textbf{RTExecutor: an Executor that eliminates nested scheduling.} We design RTExecutor, which replaces the Executor's internal callback dispatch with direct mapping to RTEMS real-time tasks. Each ROS~2 callback is assigned an independent RTEMS task with a user-specified priority, and all scheduling is performed by the RTEMS fixed-priority preemptive scheduler.

    \item \textbf{TimeManager: hardware-interrupt-driven timer callbacks.} We design TimeManager, which uses hardware tick interrupts and the RTEMS Watchdog red-black tree mechanism to drive periodic timer callbacks, reducing jitter from tens of milliseconds on Linux to $\leqslant 2$~ticks ($2$~ms) on RTEMS.

    \item \textbf{EventManager: event-driven subscription/service callback scheduling.} EventManager transfers subscription, service, and client callback execution from the Executor thread to RTEMS real-time tasks, notified via RTEMS events and protected by priority-inheritance semaphores.

    \item \textbf{Formal response-time analysis.} Since nested scheduling is eliminated, classic fixed-priority preemptive scheduling response-time analysis becomes directly applicable. We derive the schedulability conditions and show that timer release jitter is bounded by $2$~ticks.

    \item \textbf{Experimental evaluation} on the ARM QEMU platform demonstrates: (i) timer callback jitter $\leqslant 2$~ticks; (ii) deterministic priority preemption; (iii) correct concurrent callback execution with priority inheritance; (iv) end-to-end subscription latency consistent with theoretical bounds.
\end{enumerate}

The remainder of this paper is organized as follows. Section~II analyzes the nested scheduling problem. Section~III describes the ROS~2-to-RTEMS porting. Section~IV presents the RTExecutor design. Section~V provides the response-time analysis. Section~VI evaluates the system. Section~VII discusses related work. Section~VIII concludes.

\section{The Nested Scheduling Problem in ROS~2}

This section analyzes the root cause of real-time deficiencies in the standard ROS~2 execution model.
We first describe the Executor architecture, then formalise the nested scheduling problem, validate the resulting limitations through targeted motivating experiments, argue why it can be eliminated on a true RTOS, and finally present the RTEMS timer mechanism that underpins our TimeManager design.

\subsection{ROS~2 Executor Architecture}

The ROS~2 middleware dispatches user callbacks through the \texttt{rclcpp} Executor, which implements a three-phase event loop:
\begin{enumerate}
    \item \textbf{Wait.} The Executor calls \texttt{rcl\_wait}, which blocks until at least one callback-generating entity (timer, subscription, service, client) becomes ready.
    \item \textbf{Collect.} \texttt{get\_next\_executable} scans the wait-set and returns a ready callback.
    \item \textbf{Execute.} The callback function is invoked on the calling thread.
\end{enumerate}
This cycle repeats for the lifetime of the Executor.

The \texttt{SingleThreadedExecutor} runs the entire cycle on a single thread.
All ready callbacks are dispatched in FIFO order: the first callback to become ready is the first to execute, regardless of any logical priority the application may assign.
Since only one callback runs at a time, there is no concurrency within the Executor, and high-priority callbacks must wait for any lower-priority callback that entered the ready set earlier.

The \texttt{MultiThreadedExecutor} maintains a thread pool and a mutex-protected wait-set.
Multiple threads may execute callbacks concurrently, but mutual exclusion is enforced by the \emph{CallbackGroup} mechanism.
A \texttt{MutuallyExclusive} callback group ensures that at most one callback from the group executes at any time---internally implemented via a mutex.
A \texttt{Reentrant} callback group imposes no such constraint, allowing multiple callbacks from the same group to execute in parallel.
While the \texttt{MultiThreadedExecutor} introduces concurrency, it does not introduce priority-aware dispatch: ready callbacks are still assigned to available threads in FIFO or round-robin fashion.

Because the Executor implements its own scheduling logic on top of the OS thread scheduler, a \emph{nested scheduling} structure emerges: the Executor layer decides \emph{which} callback to run next, while the OS layer decides \emph{which} thread to run next.
These two layers operate independently and without priority coordination, forming the root cause of the problems we analyse below.

\subsection{Formalisation of the Nested Scheduling Problem}

We recall the classic real-time scheduling model~\cite{liu1973scheduling}.
A task set $\mathcal{T} = \{\tau_1, \tau_2, \dots, \tau_n\}$ consists of $n$ sporadic or periodic tasks, each characterized by a period (or minimum inter-arrival time) $T_i$, a worst-case execution time $C_i$, and a relative deadline $D_i$.
Under fixed-priority preemptive (FPP) scheduling, each task is assigned a static priority $\pi_i$, and the highest-priority ready task always preempts lower-priority ones.
This model provides well-established schedulability tests~\cite{liu1973scheduling} and response-time analysis.

In ROS~2, the schedulable entities are not OS tasks but \emph{callbacks}.
Let $\Gamma = \{\gamma_1, \gamma_2, \dots, \gamma_n\}$ denote the set of callbacks registered with an Executor.
Each callback $\gamma_i$ has an implicit urgency: a timer callback controlling a 1\,kHz control loop is more time-critical than a timer callback logging diagnostics at 10\,Hz.
However, the standard Executor \emph{does not model or use this priority information}.
Instead, it maintains a ready queue and dispatches callbacks in FIFO order (for \texttt{SingleThreadedExecutor}) or round-robin among threads (for \texttt{MultiThreadedExecutor}).

The nested scheduling structure in ROS~2 can be modeled as two independent schedulers operating at different abstraction layers:
\begin{itemize}
    \item \textbf{Executor layer.} A callback-level scheduler that selects the next ready callback from the wait-set using FIFO or round-robin policy. It is unaware of OS thread priorities and treats all callbacks as having equal priority.
    \item \textbf{OS layer.} A thread-level scheduler (e.g., Linux CFS or \texttt{SCHED\_FIFO}) that selects the next runnable thread based on thread priorities. It is unaware of callback priorities, since multiple callbacks of different urgencies may share the same Executor thread.
\end{itemize}

The fundamental issue is an \emph{information gap}: the Executor does not know the OS-level thread priorities, and the OS does not know the application-level callback priorities.
Consequently, the OS cannot make informed preemption decisions at the callback granularity, and the Executor cannot leverage OS-level priority mechanisms.
Formally, let $\pi^{\text{OS}}_k$ denote the OS priority of the $k$-th thread, and let $\pi^{\text{app}}_i$ denote the application-level priority of callback $\gamma_i$.
If callbacks $\gamma_i$ and $\gamma_j$ with $\pi^{\text{app}}_i > \pi^{\text{app}}_j$ are both dispatched to the same thread $k$, then the OS priority $\pi^{\text{OS}}_k$ cannot simultaneously reflect both $\pi^{\text{app}}_i$ and $\pi^{\text{app}}_j$.
The effective priority of each callback is thus lost at the Executor--OS boundary.

\subsection{Experimental Validation of Executor Scheduling Limitations}

To validate that the deficiencies above are caused by the Executor scheduling structure rather than by incidental system noise, we design a set of phase-controlled motivating experiments on RTEMS~6.1 running on the ARM \texttt{realview\_pbx\_a9\_qemu} platform (1~tick = 1~ms). All timestamps are recorded using \texttt{rtems\_clock\_get\_uptime} with sub-tick (nanosecond) resolution. The key idea is to make callback readiness observable and controllable: a low-urgency callback $\gamma_\ell$ is deliberately released slightly before a high-urgency callback $\gamma_h$, and we then measure whether $\gamma_h$ can preempt or bypass $\gamma_\ell$. Each experiment includes a FIFO-based Executor scenario (A) and an RTOS fixed-priority preemptive scenario (B) for comparison.

For each callback $\gamma_i$, we record three timestamps: its readiness time $r_i$, actual execution start time $s_i$, and finish time $f_i$. The main metric is the dispatch latency $L_i=s_i-r_i$, and the response time is $R_i=f_i-r_i$. For the high-priority callback, a dispatch latency $L_h$ that grows with the execution time of a lower-priority callback is direct evidence of callback-level priority inversion. Unless otherwise stated, callback execution time is generated by CPU-bound busy waiting rather than sleeping, so that the experiment models non-preemptive callback execution inside the Executor.

\textbf{E1: Priority inversion under varying high-priority period and task count.}
This experiment jointly evaluates two factors that determine priority inversion severity: (i)~the high-priority period $T_h$, which controls how often $\gamma_h$ encounters a busy Executor, and (ii)~the total task count $n = 1 + k$, where $k$ low-priority callbacks ($C_\ell = 10$~ms, period $= 20$~ms) compete with one high-priority callback ($C_h = 2$~ms, period $T_h$) in a \texttt{SingleThreadedExecutor}. We sweep $T_h \in \{5, 10, 20, 40\}$~ms and $n \in \{2, 3, 4, 5, 6\}$, measuring the average high-priority dispatch latency $\overline{L_h}$.

\begin{table}[t]
\centering
\caption{E1: Average high-priority dispatch latency $\overline{L_h}$ ($\mu$s) under FIFO dispatch (Scenario~A). Low-priority period = 20~ms, $C_\ell = 10$~ms, $C_h = 2$~ms.}
\label{tab:e1_fifo}
\begin{tabular}{c|ccccc}
\hline
$T_h$ (ms) \textbackslash\ $n$ & 2 & 3 & 4 & 5 & 6 \\
\hline
5   & 3\,643   & 35\,793  & 53\,531  & 65\,338  & 66\,865 \\
10  & 5\,742   & 28\,281  & 56\,399  & 64\,042  & 68\,632 \\
20  & 10\,013  & 28\,069  & 60\,048  & 73\,045  & 82\,043 \\
40  & 19       & 4\,078   & 32\,068  & 41\,063  & 75\,398 \\
\hline
\end{tabular}
\end{table}

\begin{table}[t]
\centering
\caption{E1: Average high-priority dispatch latency $\overline{L_h}$ ($\mu$s) under RTOS fixed-priority preemptive scheduling (Scenario~B). Same parameters as Table~\ref{tab:e1_fifo}.}
\label{tab:e1_rtos}
\begin{tabular}{c|ccccc}
\hline
$T_h$ (ms) \textbackslash\ $n$ & 2 & 3 & 4 & 5 & 6 \\
\hline
5   & 15.2 & 10.0 & 8.7 & 9.7 & 8.7 \\
10  & 14.5 & 18.2 & 12.4 & 14.8 & 13.8 \\
20  & 11.9 & 11.9 & 9.3 & 9.3 & 9.5 \\
40  & 15.1 & 20.2 & 22.5 & 20.3 & 22.3 \\
\hline
\end{tabular}
\end{table}

Tables~\ref{tab:e1_fifo} and~\ref{tab:e1_rtos} present the results. Under FIFO dispatch (Table~\ref{tab:e1_fifo}), $\overline{L_h}$ increases sharply with $n$: when $n=2$ (one low-priority callback), the latency is modest ($3.6$--$10$~ms depending on $T_h$), but with $n=6$ (five low-priority callbacks), it reaches $66$--$82$~ms, reflecting cumulative queueing of low-priority callbacks ahead of $\gamma_h$. The CPU utilisation of the $k$ low-priority tasks is $k C_\ell / 20\,\text{ms}$; at $k=2$ the system is fully loaded (100\%), and beyond that the FIFO queue grows without bound during each period. Under RTOS fixed-priority scheduling (Table~\ref{tab:e1_rtos}), $\overline{L_h}$ remains below $23\,\mu$s in all configurations---three to four orders of magnitude lower---because the high-priority task immediately preempts any executing low-priority task.

\textbf{E2: Ineffectiveness of OS-level real-time priority for intra-Executor inversion.}
To test whether raising the Executor thread priority can remove the inversion, we run a single low-priority and a single high-priority callback ($C_\ell = 30$~ms, $C_h = 2$~ms, $\Delta = 1$~ms, period $= 60$~ms) while sweeping the Executor thread's RTOS priority over $\{5, 10, 20, 40\}$.

\begin{table}[t]
\centering
\caption{E2: Average high-priority dispatch latency ($\mu$s) at different Executor priorities. $C_\ell = 30$~ms, $C_h = 2$~ms.}
\label{tab:e2}
\begin{tabular}{lcc}
\hline
\textbf{Setting} & $\overline{L_h}$ ($\mu$s) & $\overline{L_h}^{\max}$ ($\mu$s) \\
\hline
FIFO, priority 5  & 30\,060 & 30\,200 \\
FIFO, priority 10 & 29\,080 & 29\,142 \\
FIFO, priority 20 & 29\,090 & 29\,157 \\
FIFO, priority 40 & 29\,062 & 29\,083 \\
RTOS (separate tasks) & 16.5 & --- \\
\hline
\end{tabular}
\end{table}

Table~\ref{tab:e2} shows that $\overline{L_h}$ remains at approximately $29$~ms (i.e., $C_\ell - \Delta$) regardless of the Executor's OS priority. Raising the Executor thread to the highest RTOS priority (5) does not reduce $L_h$, because the blocking is internal to the Executor: once a low-priority callback begins executing on the Executor thread, no OS-level priority mechanism can preempt it at the callback granularity. Only when each callback runs in its own RTOS task (bottom row) does the latency drop to near zero.

\textbf{E3: Serialization caused by MutuallyExclusive callback groups.}
We evaluate the \texttt{MultiThreadedExecutor} with both $\gamma_h$ and $\gamma_\ell$ in the same \texttt{MutuallyExclusive} callback group (simulated by a shared binary semaphore). $C_\ell = 30$~ms, $C_h = 2$~ms, $\Delta = 1$~ms, period $= 60$~ms. The thread-pool size $m$ is swept over $\{1, 2, 4\}$.

\begin{table}[t]
\centering
\caption{E3: High-priority dispatch latency ($\mu$s) under MutuallyExclusive serialization.}
\label{tab:e3}
\begin{tabular}{lcc}
\hline
\textbf{Setting} & $\overline{L_h}$ ($\mu$s) & $\overline{L_h}^{\max}$ ($\mu$s) \\
\hline
Mutex, $m=1$ & 29\,338 & 30\,195 \\
Mutex, $m=2$ & 29\,076 & 29\,098 \\
Mutex, $m=4$ & 29\,059 & 29\,094 \\
RTOS (no mutex) & 16.8 & 26.7 \\
\hline
\end{tabular}
\end{table}

Table~\ref{tab:e3} confirms that increasing the thread-pool size from 1 to 4 has negligible effect: the mutual-exclusion constraint ensures that only one callback executes at a time regardless of available threads, and $\overline{L_h} \approx 29$~ms in all cases. The RTOS comparison (without the mutex, each callback in a separate prioritized task) achieves $\overline{L_h} = 16.8\,\mu$s.

\textbf{E4: Timer dispatch jitter caused by the Executor event loop.}
We evaluate whether a high-frequency timer can provide stable dispatch under callback interference. A 100~Hz high-priority timer ($C_h = 1$~ms) runs together with an occasional long-running low-priority callback ($C_\ell = 30$~ms, one-shot at $t = 50$~ms). For the $j$-th timer release, we compute the dispatch delay $s_j - e_j$ and inter-arrival jitter $|(s_j - s_{j-1}) - T|$.

\begin{table}[t]
\centering
\caption{E4: Timer dispatch delay and inter-arrival jitter ($\mu$s) with and without low-priority interference. $T = 10$~ms.}
\label{tab:e4}
\begin{tabular}{lcccc}
\hline
\textbf{Scenario} & Delay avg & Delay max & Jitter avg & Jitter max \\
\hline
FIFO Executor  & 334 & 8\,048 & 2\,261 & 28\,000 \\
RTOS priority  & 0   & 0      & 10    & 50 \\
\hline
\end{tabular}
\end{table}

Table~\ref{tab:e4} shows that under the FIFO Executor, the low-priority callback creates a dispatch delay spike of up to $8.0$~ms and a jitter spike of $28$~ms when it occupies the Executor during the timer's scheduled firing window. Under RTOS priority scheduling, both the delay and jitter remain negligible, confirming that timer dispatch is deterministic when each timer callback executes in a dedicated prioritized task.

Together, these experiments provide direct empirical evidence of the nested scheduling problem. E1 shows that FIFO dispatch latency grows with both the number of low-priority tasks and their cumulative load, E2 demonstrates that raising the OS thread priority cannot fix intra-Executor inversion, E3 confirms that \texttt{MutuallyExclusive} groups serialize regardless of thread count, and E4 shows that timer dispatch jitter spikes when the Executor is occupied. These observations motivate replacing the Executor-centric dispatch path with an RTOS-backed design in which each callback or timer service is mapped to a schedulable task with an explicit priority.

\subsection{Why Nested Scheduling Can Be Eliminated on RTOS}

A real-time operating system provides the scheduling primitives that the ROS~2 Executor lacks:
\begin{itemize}
    \item \textbf{Deterministic fixed-priority preemptive scheduling.} The RTOS kernel guarantees that the highest-priority ready task always runs, with a bounded preemption latency.
    \item \textbf{Priority inheritance protocol.} When a high-priority task is blocked by a lower-priority task holding a shared resource, the lower-priority task temporarily inherits the high priority, bounding the inversion duration.
    \item \textbf{Predictable interrupt latency.} The time from a hardware interrupt to the start of the associated ISR is bounded and typically on the order of microseconds, not milliseconds as on general-purpose Linux.
    \item \textbf{Fine-grained task priority configuration.} Each RTOS task can be assigned a distinct priority from a wide range (e.g., 1--255 in RTEMS), enabling precise differentiation of callback urgencies.
\end{itemize}

The key insight is that \emph{if each ROS~2 callback maps directly to an RTOS task with its own priority, the middleware-level scheduling performed by the Executor becomes entirely redundant}.
The RTOS scheduler already provides the priority-aware, preemptive dispatch that the Executor fails to implement.
By eliminating the Executor's internal dispatch and delegating all scheduling decisions to the RTOS kernel, we collapse the two-layer nested scheduling into a single, well-understood scheduling layer.

This approach is consistent in spirit with the CallbackIsolatedExecutor~\cite{ishikawa2025rtas}, which maps each callback to a dedicated \texttt{SCHED\_FIFO} thread on Linux.
However, CallbackIsolatedExecutor remains subject to the nondeterminism of the Linux kernel: CFS time slicing, unbounded interrupt latency, and the lack of guaranteed priority inheritance for user-space synchronization primitives.
By contrast, our approach operates on a true RTOS (RTEMS), where the scheduler provides deterministic guarantees, interrupt latency is bounded, and priority inheritance is a kernel-level primitive.
The EventsExecutor~\cite{teper2024rtss} similarly bridges ROS~2 with classic fixed-priority scheduling, but it too is constrained by the Linux kernel's nondeterministic behavior.
Our work eliminates the middleware scheduling layer entirely on an RTOS where the OS scheduler can fully and reliably take over.

\subsection{RTEMS~6.1 Timer Mechanism}

Our TimeManager design (Section~IV) leverages the native RTEMS timer mechanism.
This subsection describes the relevant internals of RTEMS~6.1.

RTEMS manages timers through two key data structures: \texttt{Timer\_Control} and \texttt{Watchdog\_Control}.
When an application creates a timer via \texttt{rtems\_timer\_create}, the kernel allocates a \texttt{Timer\_Control} block.
Firing a timer with \texttt{rtems\_timer\_server\_fire\_after} inserts a \texttt{Watchdog\_Control} node into a red-black tree indexed by expiration tick.
The red-black tree ensures $O(\log n)$ insertion and removal of timer entries, where $n$ is the number of active timers.

The execution flow from hardware interrupt to callback invocation proceeds as follows:
\begin{enumerate}
    \item \textbf{Hardware tick interrupt.} A hardware timer generates a periodic tick interrupt at the configured frequency (e.g., 1\,kHz).
    \item \textbf{Clock Tick ISR.} The interrupt service routine increments the system tick counter and triggers the Watchdog chain.
    \item \textbf{Watchdog expiration.} The kernel traverses the red-black tree and identifies all \texttt{Watchdog\_Control} nodes whose expiration tick is less than or equal to the current tick. Each expired watchdog is removed from the tree and enqueued for processing.
    \item \textbf{Timer Server execution.} Expired watchdogs are dispatched to the \emph{Timer Server} task---a dedicated RTEMS task that executes the registered callback functions in task context (not ISR context), ensuring that callbacks can invoke blocking OS services.
\end{enumerate}

A critical feature of this mechanism is the priority ordering of the involved tasks.
In our configuration, the Timer Server task runs at the highest application priority (priority~1 in RTEMS, where lower numbers indicate higher priority), followed by Worker tasks at priority~5, and the Executor spin thread at priority~50.
This strict priority hierarchy ensures that:
\begin{itemize}
    \item The Timer Server preempts any Worker or Executor task, guaranteeing that timer callbacks are dispatched with minimal latency after the hardware tick.
    \item Worker tasks executing subscription or service callbacks preempt the spin thread, preventing readiness detection from interfering with callback execution.
    \item Priority inheritance applies when tasks contend for shared resources (e.g., the callback ready queue), bounding any remaining priority inversion.
\end{itemize}

This hardware-driven, priority-ordered timer mechanism provides the foundation for our TimeManager design.
By registering ROS~2 timer callbacks as RTEMS timer server callbacks, we achieve jitter bounded by the tick resolution (1\,ms in our configuration), a dramatic improvement over the software timer approach in standard \texttt{rclcpp}, where timer expiration is checked only when the Executor's event loop reaches the wait phase.
\section{Porting ROS~2 to RTEMS}
\label{sec:porting}

This section describes the complete porting of the ROS~2 software stack to the RTEMS~6.1 real-time operating system on the ARM platform.
Unlike micro-ROS, which targets constrained microcontrollers by providing only a subset of the ROS~2 API, our port preserves the full \texttt{rclcpp} programming model---including intra-process communication, lifecycle nodes, and the complete message ecosystem---so that existing ROS~2 applications can be rebuilt for RTEMS with minimal source changes.
We organize the discussion around the five major challenges encountered during porting: platform selection and scope (Section~\ref{subsec:overview}), POSIX compatibility (Section~\ref{subsec:posix}), \texttt{rclcpp} adaptation (Section~\ref{subsec:rclcpp_adapter}), DDS communication (Section~\ref{subsec:dds}), and build/deployment (Section~\ref{subsec:build}).

\subsection{Porting Overview}
\label{subsec:overview}

The target platform is the ARM \texttt{realview\_pbx\_a9\_qemu} Board Support Package (BSP) under RTEMS~6.1.
This BSP provides a dual-core ARM Cortex-A9 emulation environment with a Generic Interrupt Controller (GIC), a hardware timer, and a VirtIO network device, making it suitable for both functional validation and real-time evaluation without physical hardware.

Table~\ref{tab:ported_components} summarizes the ROS~2 components ported.
The scope covers the entire \texttt{rclcpp} stack: the \texttt{rcl} C library, the \texttt{rclcpp} C++ library, the \texttt{rmw} abstraction layer, the \texttt{rmw\_fastrtps\_cpp} implementation, the \texttt{rosidl} message generation pipeline, the eProsima Fast-DDS middleware, and all standard message and service packages required by the \texttt{rclcpp} API.
In total, 26 ROS~2 packages were cross-compiled and statically linked for the RTEMS target.

\begin{table}[t]
\centering
\caption{Ported ROS~2 components}
\label{tab:ported_components}
\begin{tabular}{ll}
\hline
\textbf{Layer} & \textbf{Components} \\
\hline
Middleware    & Fast-DDS, rmw, rmw\_fastrtps\_cpp \\
Core C library & rcl \\
Core C++ library & rclcpp \\
IDL / Messages & rosidl (adapter, typesupport, \\
               & runtime), std\_msgs, geometry\_msgs, \\
               & sensor\_msgs, \ldots \\
Build tools   & ament\_cmake, RTcolcon \\
\hline
\end{tabular}
\end{table}

A key design decision was to port the \emph{full} \texttt{rclcpp} stack rather than the micro-ROS subset.
Micro-ROS uses the Micro XRCE-DDS agent/client architecture, which introduces an additional network hop and does not support intra-process zero-copy communication---a feature critical for low-latency callback chains on a single processor.
By retaining the standard \texttt{rclcpp} intra-process path, our port enables the RTExecutor (Section~\ref{sec:design}) to exploit shared-memory zero-copy data transfer within the same address space, avoiding serialization overhead entirely.

The build system is based on cross-compilation: all packages are compiled on an x86-64 host using the \texttt{arm-rtems6-gcc}/\texttt{arm-rtems6-g++} cross-compiler with RTEMS~6.1 BSP headers and libraries.
A custom build orchestrator, RTcolcon (Section~\ref{subsec:build}), manages the 26-component dependency graph and parallel compilation.

\subsection{POSIX Compatibility Layer}
\label{subsec:posix}

RTEMS implements a POSIX-compliant subset (POSIX~1003.1-2008 profile), which covers most of the interfaces used by ROS~2's C libraries.
However, several Linux-specific or GNU-extension APIs that \texttt{rcl} and its dependencies rely on are absent.
Table~\ref{tab:posix_adaptations} summarizes the key adaptations.

\begin{table}[t]
\centering
\caption{POSIX API adaptations for RTEMS}
\label{tab:posix_adaptations}
\begin{tabular}{p{2.1cm}p{2.1cm}p{3.0cm}}
\hline
\textbf{Linux API} & \textbf{RTEMS Replacement} & \textbf{Notes} \\
\hline
\texttt{pthread\_create} & \texttt{rtems\_task\_create} / \texttt{rtems\_task\_start} & RTEMS Classic API tasks \\
\texttt{pthread\_mutex} & RTEMS binary semaphore & \texttt{RTEMS\_INHERIT\_PRIORITY} attribute \\
\texttt{pthread\_cond} & RTEMS event + semaphore & Signal via \texttt{rtems\_event\_send}; wait via \texttt{rtems\_event\_receive} \\
\texttt{clock\_gettime (CLOCK\_MONOTONIC)} & \texttt{rtems\_clock\_get\_ticks\_since\_boot} & Tick-based monotonic time \\
\texttt{on\_exit} (GNU ext.) & \texttt{atexit} wrapper & Guarded by \texttt{\#ifdef RTEMS} \\
\texttt{pthread\_setaffinity\_np} & Skipped & Single-core target; no affinity needed \\
\hline
\end{tabular}
\end{table}

\textbf{Threads.}
The \texttt{rcl} and \texttt{rclcpp} libraries create internal threads for the DDS participant, the wait-set polling loop, and the executor.
Under Linux these are \texttt{pthread} objects; under RTEMS we replace them with Classic API tasks created via \texttt{rtems\_task\_create()} and started via \texttt{rtems\_task\_start()}.
This mapping preserves the ability to assign each task an independent RTEMS priority, which is essential for the RTExecutor design (Section~\ref{sec:design}).

\textbf{Mutexes.}
ROS~2 uses \texttt{pthread\_mutex} extensively to protect shared data structures such as the wait-set and callback groups.
On RTEMS, we replace each mutex with a binary semaphore created with the \texttt{RTEMS\_INHERIT\_PRIORITY} attribute.
This ensures that when a low-priority task holds a lock needed by a higher-priority task, the low-priority task's priority is temporarily elevated---a direct application of the priority inheritance protocol that prevents unbounded priority inversion at the OS level.

\textbf{Condition variables.}
RTEMS does not natively implement \texttt{pthread\_cond\_wait} / \texttt{pthread\_cond\_signal}.
We emulate condition-variable semantics using a combination of RTEMS events and semaphores: the waiting task calls \texttt{rtems\_event\_receive()} with the \texttt{RTEMS\_WAIT} option, and the signaling task calls \texttt{rtems\_event\_send()} to wake it.
A binary semaphore serializes access to the shared predicate state.
This pattern closely mirrors the standard condvar protocol (lock--check--wait loop) and introduces minimal overhead.

\textbf{Monotonic clock.}
The \texttt{rcl} timer subsystem relies on \texttt{clock\_gettime(CLOCK\_MONOTONIC)} to compute timeouts and periods.
RTEMS provides \texttt{rtems\_clock\_get\_ticks\_since\_boot()}, which returns a monotonically non-decreasing tick count.
We wrap this call behind a thin compatibility function that converts ticks to nanoseconds using the configured tick rate, preserving the \texttt{rcl} timer API without modification.

\textbf{GNU extensions.}
The \texttt{on\_exit()} function is a GNU extension not provided by RTEMS.
We replace it with an \texttt{atexit()}-based wrapper guarded by \texttt{\#ifdef RTEMS}.
Similarly, \texttt{pthread\_setaffinity\_np()} is a Linux-specific call for setting CPU affinity.
Since our target platform runs a single-core configuration, we skip this call entirely under the same guard.

All adaptations are confined to a small set of header files and do not alter the public API of \texttt{rcl} or \texttt{rclcpp}, ensuring that application code remains source-compatible across Linux and RTEMS.

\subsection{rclcpp Adapter Layer}
\label{subsec:rclcpp_adapter}

Above the POSIX compatibility layer, \texttt{rclcpp} uses C++ standard-library abstractions---\texttt{std::thread}, \texttt{std::mutex}, \texttt{std::condition\_variable}, and \texttt{std::chrono}---that must be mapped to RTEMS primitives.
We implement an adapter layer that provides these mappings transparently.

\textbf{Thread abstraction.}
The \texttt{std::thread} constructor is intercepted so that the callable and its arguments are packaged into a structure and passed to \texttt{rtems\_task\_start()}.
The resulting RTEMS task receives the caller's requested priority through a thread-local attribute stored before task creation.
This allows the RTExecutor to create worker tasks with specific priorities directly from C++ code without exposing RTEMS-specific APIs to the application.

\textbf{Mutex.}
Each \texttt{std::mutex} object is backed by an RTEMS binary semaphore with the \texttt{RTEMS\_INHERIT\_PRIORITY} attribute.
The \texttt{lock()} call invokes \texttt{rtems\_semaphore\_obtain()} with \texttt{RTEMS\_WAIT}, and \texttt{unlock()} invokes \texttt{rtems\_semaphore\_release()}.
Because priority inheritance is enforced at the kernel level, any priority inversion between \texttt{rclcpp} threads is bounded automatically.

\textbf{Condition variable.}
\texttt{std::condition\_variable::wait()} is implemented by releasing the associated mutex semaphore, calling \texttt{rtems\_event\_receive()} on a per-condition-variable event set, and re-acquiring the mutex upon wake-up.
\texttt{notify\_one()} and \texttt{notify\_all()} map to \texttt{rtems\_event\_send()} targeting one or all waiting tasks, respectively.
Spurious wake-ups are handled by the standard predicate-loop idiom already present in \texttt{rclcpp}.

\textbf{Time source.}
\texttt{std::chrono::steady\_clock} is specialized for RTEMS by reading \texttt{rtems\_clock\_get\_ticks\_since\_boot()} and converting the result to a \texttt{time\_point} using the board's configured microseconds-per-tick constant.
All \texttt{rclcpp} time queries---timer period computation, deadline tracking, and duration measurements---then operate correctly without any changes to higher-level code.

\textbf{C++17 compatibility.}
The RTEMS cross-compiler (\texttt{arm-rtems6-g++}) supports C++17 with \texttt{-std=c++17}, but some GNU extensions used by \texttt{rclcpp} (e.g., \texttt{\_\_attribute\_\_((cleanup))}, certain thread-local storage qualifiers) are either unavailable or behave differently in strict conformance mode.
We resolve these through conditional compilation (\texttt{\#ifdef \_\_rtems\_\_}) and, where necessary, replace GNU-specific constructs with standard C++17 equivalents.
For example, GNU-style statement expressions within macros are rewritten as inline template functions.

\subsection{DDS and Communication Stack}
\label{subsec:dds}

The communication backbone of ROS~2 is the Data Distribution Service (DDS).
We selected eProsima Fast-DDS as the DDS implementation because it is the default in ROS~2 and has a relatively clean CMake-based build system amenable to cross-compilation.

\textbf{Fast-DDS porting.}
Fast-DDS relies on POSIX networking APIs (\texttt{socket}, \texttt{bind}, \texttt{sendto}, \texttt{recvfrom}) for UDP multicast discovery and TCP data transport.
The \texttt{realview\_pbx\_a9\_qemu} BSP provides a VirtIO network device with a standard BSD socket interface, which Fast-DDS uses without modification after adjusting the interface name detection logic (RTEMS uses \texttt{/dev/eth0} rather than Linux's \texttt{eth0}).
Thread creation inside Fast-DDS is routed through the POSIX compatibility layer described in Section~\ref{subsec:posix}.

\textbf{rmw\_fastrtps\_cpp.}
The \texttt{rmw\_fastrtps\_cpp} package bridges the \texttt{rmw} abstraction with Fast-DDS.
Under Linux, ROS~2 typically loads this package as a shared library at runtime via \texttt{dlopen()}.
RTEMS does not support dynamic loading; therefore, \texttt{rmw\_fastrtps\_cpp} is statically linked into the final executable.
The \texttt{rmw} initialization path is adjusted to call the Fast-DDS registration function directly rather than through the dynamic loader.

\textbf{Intra-process zero-copy communication.}
ROS~2's intra-process feature bypasses DDS entirely: when a publisher and a subscriber reside in the same process, data is transferred via \texttt{std::shared\_ptr} through a lock-free intra-process manager, avoiding serialization and network overhead.
This mechanism works natively on RTEMS because it relies solely on shared memory within the same address space---a property that holds trivially on a single-process RTOS application.
The intra-process path is particularly important for our design: it ensures that the RTExecutor's EventManager can deliver subscription data to worker tasks with minimal latency, without crossing a DDS boundary.

\textbf{Message generation.}
The \texttt{rosidl} pipeline generates C/C++ type-support code from \texttt{.msg} and \texttt{.srv} definitions.
We cross-compile the generator tools for the host and run them natively, then cross-compile the generated sources for the RTEMS target.
The generated type-support headers and Fast-DDS serialization code compile without modification once the POSIX and C++17 compatibility layers are in place.

\subsection{Build and Deployment}
\label{subsec:build}

Building ROS~2 for RTEMS requires a cross-compilation toolchain and a build orchestrator capable of handling the 26-package dependency graph.

\textbf{Cross-compilation toolchain.}
The toolchain file configures CMake to use \texttt{arm-rtems6-gcc}/\texttt{arm-rtems6-g++} as the compiler and linker, with sysroot pointing to the RTEMS~6.1 installation directory containing the BSP headers and libraries.
Compiler flags include \texttt{-mcpu=cortex-a9}, \texttt{-mfloat-abi=softfp}, and \texttt{-std=c++17}.
Linker flags enforce static linking (\texttt{-static}) and include the RTEMS BSP library (\texttt{-lrtemsbsp}).

\textbf{RTcolcon.}
The standard ROS~2 build tool, \texttt{colcon}, depends on Python~3 and host-executable CMake packages that are not available in the cross-compilation environment.
We developed RTcolcon, a custom build orchestrator that:
(i) resolves the 26-package dependency graph;
(ii) invokes CMake for each package with the cross-compilation toolchain file;
(iii) parallelizes independent builds across available host cores;
and (iv) installs headers and libraries to a staging directory consumed by downstream packages.
RTcolcon produces a single merged \texttt{CMakeCache} that ensures all packages agree on the target architecture, compiler, and RTEMS version.

\textbf{RTEMS system configuration.}
RTEMS applications are configured at compile time through the \texttt{confdefs.h} mechanism.
Key configuration parameters include:
\begin{itemize}
    \item \texttt{CONFIGURE\_UNLIMITED\_OBJECTS} --- allows the kernel to allocate tasks, semaphores, and other objects on demand rather than from a fixed pool;
    \item \texttt{CONFIGURE\_POSIX\_API} --- enables the POSIX subset required by \texttt{rcl} and Fast-DDS;
    \item \texttt{CONFIGURE\_MAXIMUM\_PROCESSORS} = 1 --- single-core configuration matching the QEMU target;
    \item \texttt{CONFIGURE\_MICROSECONDS\_PER\_TICK} = 1000 --- 1~ms tick period for timer resolution;
    \item \texttt{CONFIGURE\_INIT\_TASK\_ENTRY\_POINT} --- set to the ROS~2 application's \texttt{main()} function.
\end{itemize}

\textbf{Static linking and deployment.}
All ROS~2 libraries, Fast-DDS, and the RTEMS BSP are statically linked into a single ELF executable (\texttt{.exe}).
This executable is loaded by QEMU via the \texttt{-kernel} flag.
No filesystem or dynamic loader is required at runtime, which eliminates shared-library versioning issues and reduces boot time.
The resulting binary is self-contained: it initializes the RTEMS kernel, starts the network stack, and launches the ROS~2 application in a deterministic boot sequence.

\section{RTExecutor: Eliminating Nested Scheduling}
\label{sec:design}

\subsection{Core Principle}

The fundamental design principle of RTExecutor is that \emph{the middleware layer should not make scheduling decisions}; scheduling authority should be entirely exercised by the OS kernel.  In the standard ROS~2 execution model, the Executor maintains its own callback dispatch policy (FIFO or round-robin) inside the middleware, forming a nested scheduling structure with the OS thread scheduler.  RTExecutor eliminates this middleware-level scheduling by mapping every callback directly to an RTEMS real-time task, so that all preemption, priority arbitration, and resource sharing decisions are made by a single scheduler---the RTEMS fixed-priority preemptive kernel.

Table~\ref{tab:comparison} summarizes the key differences between the standard \texttt{rclcpp} Executor and RTExecutor.

\begin{table}[t]
\centering
\caption{Comparison of Standard Executor and RTExecutor}
\label{tab:comparison}
\begin{tabular}{p{1.3cm}|p{2.5cm}|p{2.5cm}}
\hline
\textbf{Aspect} & \textbf{Standard Executor} & \textbf{RTExecutor} \\
\hline
\hline
Callback Readiness Detection & \texttt{rcl\_wait} (middleware) & \texttt{rcl\_wait} (middleware) \\
\hline
Callback Execution Scheduling & Executor thread pool (middleware) & RTEMS tasks (OS kernel) \\
\hline
Timer Driver & \texttt{rclcpp} software timer & Hardware tick interrupt \\
\hline
Priority Management & Callbacks have no priority & Each callback mapped to RTEMS task with independent priority \\
\hline
Nested Scheduling & Exists (Executor + OS) & Eliminated (OS only) \\
\hline
\end{tabular}
\end{table}

Note that both the standard Executor and RTExecutor use \texttt{rcl\_wait} for \emph{readiness detection}---the point of divergence is what happens \emph{after} a callback becomes ready.  The standard Executor dispatches it to its own thread pool, whereas RTExecutor delegates execution to an RTEMS task, thereby transferring scheduling authority to the kernel.

\subsection{Architecture Overview}

RTExecutor inherits from \texttt{rclcpp::Executor} and overrides its core execution methods.  As shown in Fig.~\ref{fig:architecture}, RTExecutor decomposes callback management into two sub-managers:

\begin{itemize}
    \item \textbf{TimeManager} --- manages timer callbacks, driven by hardware tick interrupts rather than software timers.
    \item \textbf{EventManager} --- manages subscription, service, and client callbacks, dispatching them to RTEMS tasks via an event-driven mechanism.
\end{itemize}

The main \texttt{spin()} loop proceeds as follows:
\begin{enumerate}
    \item Start the TimeManager, which spawns a Reader--Worker task pair for hardware-timer-driven callbacks.
    \item Start the EventManager, which spawns a configurable number of Worker RTEMS tasks.
    \item Call \texttt{get\_next\_executable()} to detect a ready callback via \texttt{rcl\_wait}.
    \item Call \texttt{dispatch\_any\_executable()} to wrap the callback and push it to the appropriate manager.
\end{enumerate}

This architecture ensures that the RTExecutor spin thread performs only \emph{detection and dispatch}---it never executes callback logic itself.  All callback execution occurs within RTEMS tasks under the control of the kernel scheduler.

\begin{figure}[t]
\centering
\framebox{\parbox{0.9\columnwidth}{\centering\vspace{1em}
\textbf{RTExecutor} \texttt{(rclcpp::Executor)}\\[0.5em]
\begin{tabular}{c}
TimeManager \\ \hline
Reader Task (prio 80) \\
Worker Task (prio 70)
\end{tabular}
\quad
\begin{tabular}{c}
EventManager \\ \hline
Worker Task(s) (prio 60) \\
RtsemQueue
\end{tabular}\\[0.8em]
\texttt{spin()} $\to$ \texttt{get\_next\_executable()} $\to$ \texttt{dispatch\_any\_executable()}\\[0.3em]
Spin Thread (prio 50)
\vspace{1em}}}
\caption{RTExecutor architecture overview.  TimeManager handles timer callbacks via hardware interrupts; EventManager handles subscription/service/client callbacks via RTEMS events.  The spin thread performs only readiness detection and dispatch.}
\label{fig:architecture}
\end{figure}

\subsection{TimeManager: Hardware Interrupt-Driven Timer Callbacks}

\subsubsection{Design Basis}

The TimeManager is based on the RTEMS Timer Server mechanism, whose operation can be traced from hardware interrupt to callback execution as follows:

\begin{equation}
\text{Hardware Tick Interrupt} \to \text{Clock Tick ISR} \to \text{Watchdog R-B Tree} \to \text{Timer Server Task} \to \text{Callback}
\end{equation}

On each hardware tick, the Clock Tick ISR traverses the Watchdog red-black tree to identify expired timers.  Expired entries are posted to the Timer Server task, which executes the registered callback in its own task context.  Because the Timer Server runs at the highest application priority (priority~1), it preempts all other tasks immediately upon notification, ensuring minimal dispatch latency.

\subsubsection{Custom Timer Driver}

To expose this mechanism to ROS~2, we implement a custom RTEMS device driver, \texttt{/dev/rtss\_timer}, which provides:

\begin{itemize}
    \item Up to \textbf{8 timer channels}, each independently configurable with its own period.
    \item Three \textbf{trigger modes} selectable per channel:
    \begin{itemize}
        \item \texttt{EVENT} --- fires via \texttt{rtems\_event\_send} to the Reader task;
        \item \texttt{BINARY\_SEMAPHORE} --- releases an RTEMS binary semaphore;
        \item \texttt{MESSAGE\_QUEUE} --- enqueues a message to an RTEMS message queue.
    \end{itemize}
\end{itemize}

The driver maintains a lock-free ring buffer of 64 entries for pending timer events, enabling $O(1)$ insertion by the ISR and $O(1)$ consumption by the Reader task without lock contention.

\subsubsection{Reader--Worker Dual-Task Architecture}

The TimeManager employs a dual-task architecture to separate hardware event reading from callback execution:

\begin{itemize}
    \item \textbf{Reader task} (priority~80): blocks on \texttt{rtems\_event\_receive} waiting for hardware tick events from the timer driver.  Upon receiving an event, it identifies which timer channel fired and invokes the corresponding registered callback.
    \item \textbf{Worker task} (priority~70): executes the registered timer callbacks.  Separating the Reader from the Worker ensures that reading events from the driver is not delayed by a long-running callback, and vice versa.
\end{itemize}

\subsubsection{Registration Interface}

Timer callbacks are registered via:
\begin{center}
\texttt{register\_rt\_timer(callback, period\_ms)}
\end{center}
where \texttt{callback} is a \texttt{std::function<void()>} and \texttt{period\_ms} specifies the desired period in milliseconds.  Internally, this converts the period to ticks using the system's configured \texttt{CONFIGURE\_MICROSECONDS\_PER\_TICK} and programs the corresponding timer channel in the driver.

\subsubsection{Jitter Measurement}

The theoretical fire time for the $k$-th invocation of a timer with period $T$ (in ticks) and start tick $t_0$ is:
\begin{equation}
    t_{\text{theory}}(k) = t_0 + k \cdot T
\end{equation}
The measured jitter is:
\begin{equation}
    J(k) = t_{\text{actual}}(k) - t_{\text{theory}}(k)
\end{equation}

In our experiments (Section~VI), $J(k) \in [0,\, 2]$~ticks, i.e., $\leqslant 2$~ms with a 1~ms tick resolution.

\subsubsection{Comparison with Standard \texttt{rclcpp} Timers}

The standard \texttt{rclcpp} timer is implemented on top of \texttt{rcl\_timer}, a software timer that relies on \texttt{rcl\_wait} to detect expiration.  On Linux, this software timer exhibits typical jitter of $10$--$50$~ms due to nondeterministic scheduling, tick resolution, and CFS preemption behavior.  In contrast, the TimeManager's hardware-interrupt-driven approach achieves jitter $\leqslant 2$~ms, a reduction of more than one order of magnitude.

\subsection{EventManager: Event-Driven Subscription Callback Scheduling}

\subsubsection{Overview}

The EventManager transfers the execution of subscription, service, and client callbacks from the Executor thread to RTEMS real-time tasks.  This is the key mechanism that eliminates middleware-level scheduling for event-driven callbacks.

\subsubsection{Workflow}

The dispatch workflow for a subscription callback proceeds as follows:

\begin{enumerate}
    \item The spin thread calls \texttt{rcl\_wait}, which blocks until at least one callback is ready.
    \item Upon detecting a ready callback, the spin thread wraps it as a \texttt{std::function<void()>} and pushes it into the \texttt{RtsemQueue}.
    \item The spin thread notifies an idle Worker task via \texttt{rtems\_event\_send(RTEMS\_EVENT\_1)}.
    \item The RTEMS Worker task, which was blocked on \texttt{rtems\_event\_receive}, wakes up, retrieves the callback from the queue via \texttt{try\_pop}, and executes it.
\end{enumerate}

By transferring callback execution from the spin thread (which has no application-level priority semantics) to a dedicated RTEMS task with a user-specified priority, the EventManager ensures that all scheduling decisions are made by the RTEMS kernel.

\subsubsection{RtsemQueue Design}

The \texttt{RtsemQueue} is a task-safe queue protected by an RTEMS priority-inheritance binary semaphore.  Its two core operations are:

\begin{itemize}
    \item \textbf{\texttt{push}}: acquires the semaphore in blocking mode (\texttt{RTEMS\_WAIT}), enqueues the callback, and releases the semaphore.  This guarantees that the push operation eventually succeeds even under contention.
    \item \textbf{\texttt{try\_pop}}: attempts to acquire the semaphore in non-blocking mode (\texttt{RTEMS\_NO\_WAIT}).  If the semaphore is held by another task, \texttt{try\_pop} returns immediately without blocking, preventing the Worker task from stalling on an empty or contended queue.
\end{itemize}

The priority-inheritance attribute (\texttt{RTEMS\_INHERIT\_PRIORITY}) on the semaphore ensures that if a high-priority Worker is blocked waiting for the semaphore held by a lower-priority task, the holder's priority is temporarily elevated, bounding priority inversion to at most one critical section duration.

\subsubsection{Configurable Worker Tasks}

The EventManager supports a configurable number of Worker tasks, each created with an independently specified priority via \texttt{rtems\_task\_create}.  This allows the application developer to map different callbacks to Workers of different priorities, enabling fine-grained priority-based preemption among callbacks---a capability that is impossible under the standard Executor's uniform thread pool.

\subsection{Dispatch Mechanism: \texttt{dispatch\_any\_executable}}

After \texttt{get\_next\_executable()} detects a ready callback, \texttt{dispatch\_any\_executable()} wraps it as a lambda and pushes it to the EventManager.  The wrapping logic depends on the callback type:

\begin{itemize}
    \item \textbf{Timer callback}: wrapped as
    \begin{center}
    \texttt{[\,timer\,]() \{ timer->execute\_callback(); \}}
    \end{center}
    \item \textbf{Subscription callback}: wrapped as
    \begin{center}
    \texttt{[\,subscription\,]() \{ subscription->take(); handle(); \}}
    \end{center}
\end{itemize}

In both cases, the lambda captures a \texttt{shared\_ptr} to the callback object by value.  This is critical: because the callback is executed asynchronously by an RTEMS task rather than by the spin thread, the \texttt{shared\_ptr} capture ensures that the callback object remains alive throughout the RTEMS task's execution, preventing use-after-free errors even if the spin thread proceeds to destroy or replace the object.

For callbacks belonging to a \texttt{MutuallyExclusive} callback group, RTExecutor resets the group's \texttt{can\_be\_taken\_from} flag only \emph{after} the RTEMS task completes execution.  This preserves the mutual exclusion semantics of the callback group across the asynchronous dispatch boundary.

\subsection{Priority Mapping Strategy}

RTExecutor assigns distinct RTEMS priorities to each system task, creating a clear priority hierarchy that reflects the timing requirements of each component.  Table~\ref{tab:priority} presents the default priority mapping.

\begin{table}[t]
\centering
\caption{Default Priority Mapping in RTExecutor}
\label{tab:priority}
\begin{tabular}{l|c|l}
\hline
\textbf{RTEMS Task} & \textbf{Priority} & \textbf{Role} \\
\hline
\hline
Timer Server & 1 & Hardware interrupt callback dispatch \\
\hline
TimeManager Reader & 80 & Read hardware tick events \\
\hline
TimeManager Worker & 70 & Execute timer callbacks \\
\hline
EventManager Worker & 60 & Execute subscription/service callbacks \\
\hline
RTExecutor spin thread & 50 & Callback readiness detection and dispatch \\
\hline
Non-real-time tasks & 200 & Logging and debugging \\
\hline
\end{tabular}
\end{table}

In RTEMS, lower numeric values denote higher scheduling priority.  The Timer Server at priority~1 thus preempts all other application tasks, ensuring that timer expirations are processed with minimal latency.  The TimeManager Reader (priority~80) is given a higher priority than the Worker (priority~70) so that reading events from the hardware driver is never delayed by a running callback.

\subsubsection{Per-Callback Priority Assignment}

Beyond the default mapping, each individual callback can be assigned to an independent-priority RTEMS task.  This is achieved by extending the EventManager's \texttt{num\_workers} and configuring each Worker with a user-specified priority during registration.  Consequently, a high-priority control callback can preempt a lower-priority logging callback even though both are subscription callbacks---a distinction that the standard Executor cannot make.

\subsubsection{Priority Inheritance Support}

RTExecutor leverages the RTEMS semaphore's \texttt{RTEMS\_INHERIT\_PRIORITY} attribute to implement priority inheritance.  When a high-priority Worker is blocked on the \texttt{RtsemQueue} semaphore held by a lower-priority Worker, the lower-priority task's priority is temporarily elevated to that of the blocked task.  This bounds the worst-case blocking time $B_i$ for any task $\tau_i$ to the duration of at most one critical section of any lower-priority task that shares a protected resource.

\subsection{Implementation Details}

The RTExecutor implementation spans several source files, each with a clearly defined responsibility.

\subsubsection{RTExecutor and EventManager}

\texttt{rt\_executor.hpp} (located in \texttt{rclcpp/include/rclcpp/executors/}) declares the \texttt{RTExecutor} class and the \texttt{EventManager} class.  \texttt{rt\_executor.cpp} (in \texttt{rclcpp/src/rclcpp/executors/}) provides the complete implementation, including the \texttt{TimeManagerImpl} class that encapsulates all TimeManager logic.

\subsubsection{Hardware Timer Driver}

\texttt{rtss\_timer\_driver.c} (in \texttt{intra\_process\_demo/driver/}) implements the custom RTEMS device driver.  The driver uses a lock-free ring buffer of 64 entries for ISR-to-task communication, avoiding lock contention between the interrupt context and task context.

\subsubsection{RtsemQueue}

The \texttt{RtsemQueue} is a task-safe queue built on top of the RTEMS priority-inheritance binary semaphore.  Its \texttt{push} and \texttt{try\_pop} operations provide the synchronization interface between the spin thread (producer) and the Worker tasks (consumers).

\subsubsection{EventManager Worker Task Loop}

Each EventManager Worker task executes the following loop:
\begin{enumerate}
    \item Block on \texttt{rtems\_event\_receive(RTEMS\_EVENT\_1, RTEMS\_WAIT, ...)}.
    \item Upon receiving the event, call \texttt{RtsemQueue::try\_pop()} to retrieve a callback.
    \item If a callback is obtained, execute it.
    \item Return to step~1.
\end{enumerate}
This event-driven design ensures that Worker tasks consume no CPU time while idle, and wake up with minimal latency when a callback is dispatched.

\subsubsection{TimeManagerImpl Reader--Worker Architecture}

The \texttt{TimeManagerImpl} class internally creates two RTEMS tasks:
\begin{itemize}
    \item The \textbf{Reader task} blocks on a timer driver event, reads the fired channel information, and signals the Worker task.
    \item The \textbf{Worker task} executes the registered timer callback corresponding to the fired channel.
\end{itemize}
The separation ensures that ISR event processing is never delayed by a long-running callback.

\subsubsection{Lifecycle Management in \texttt{dispatch\_any\_executable}}

The \texttt{shared\_ptr} capture in the dispatch lambda is the cornerstone of safe asynchronous callback execution.  When \texttt{dispatch\_any\_executable()} pushes a lambda into the \texttt{RtsemQueue}, the captured \texttt{shared\_ptr} increments the callback object's reference count.  The count is decremented only when the RTEMS Worker task finishes executing the lambda and the lambda object is destroyed.  This guarantees that the callback object's lifetime extends across the full duration of its execution in the RTEMS task context, even if the spin thread has already moved on to process subsequent callbacks.

\subsubsection{Intra-Process Callback Execution Flow}

When intra-process (zero-copy) communication is enabled, the publisher's callback writes data to shared memory, and the subscription callback reads it from the same memory region without serialization.  In RTExecutor, both the publisher and subscription callbacks execute in their respective RTEMS tasks.  The EventManager dispatches the subscription callback to a Worker task upon detecting readiness via \texttt{rcl\_wait}.  The zero-copy data path and the RTEMS task-based execution are orthogonal and compose seamlessly: the shared memory is accessed under the protection of the \texttt{RtsemQueue} semaphore with priority inheritance, ensuring that data races are prevented and priority inversion is bounded.

\section{Response-Time Analysis}

This section derives the response-time bounds for callbacks managed by RTExecutor.
Since the nested scheduling structure is eliminated---all callbacks execute as RTEMS tasks under a single fixed-priority preemptive scheduler---classic real-time scheduling theory becomes directly applicable.

\subsection{Scheduling Model}

We adopt the following assumptions:
\begin{enumerate}
    \item The RTEMS kernel employs fixed-priority preemptive scheduling with FIFO tie-breaking (FP-FIFO) on a single-core processor.
    \item Each ROS~2 callback $\gamma_i$ in the callback set $\Gamma = \{\gamma_1, \gamma_2, \dots, \gamma_n\}$ is mapped to a dedicated RTEMS task $\tau_i$ with priority $\pi_i$, period $T_i$, worst-case execution time $C_i$, and relative deadline $D_i$.
    \item Priorities are distinct and assigned according to the rate-monotonic policy: a shorter period implies a higher priority.
    \item RTEMS priority inheritance is enabled on all shared resources (e.g., the \texttt{RtsemQueue}), bounding priority inversion.
\end{enumerate}

Because RTExecutor delegates all scheduling decisions to the RTEMS kernel, the middleware-layer scheduler no longer exists as an independent scheduling entity.
Consequently, the response-time analysis degenerates to the well-studied classic fixed-priority preemptive scheduling model, and no Executor-level blocking terms need to be accounted for.

\subsection{Timer Callback Response Time}

\subsubsection{Release Model}

The TimeManager releases timer callbacks through the following chain:
\begin{equation}
    \text{HW tick interrupt} \;\to\; \text{Watchdog expiration} \;\to\; \text{Timer Server notification} \;\to\; \text{Worker execution}.
\end{equation}
At each clock tick, the Clock Tick ISR fires and inspects the RTEMS Watchdog red-black tree.
If a timer's expiration time has been reached, the Timer Server task is notified, which in turn signals the corresponding Worker task via an RTEMS event or semaphore.
The Worker task then executes the registered timer callback.

\subsubsection{Release Jitter}

The timer callback release suffers from jitter $J_i^{\mathrm{timer}}$ bounded by:
\begin{equation}
    J_i^{\mathrm{timer}} \leqslant 2 \; \text{ticks}.
\end{equation}
This bound arises from four sources:
\begin{itemize}
    \item \textbf{Tick quantization error} ($\leqslant 1$~tick): The hardware timer fires at the next tick boundary after the requested expiration time, which may be up to one tick early or late relative to the ideal release instant.
    \item \textbf{ISR execution delay}: The Clock Tick ISR must traverse the Watchdog red-black tree; the traversal time adds a bounded delay before the Timer Server is notified.
    \item \textbf{Timer Server context switch}: The Timer Server task must be dispatched by the scheduler before it can signal the Worker task.
    \item \textbf{Worker notification}: The \texttt{rtems\_event\_send} operation and the subsequent context switch to the Worker task contribute at most one additional tick of delay.
\end{itemize}
In our experimental platform (ARM \texttt{realview\_pbx\_a9\_qemu} with $\text{CONFIGURE\_MICROSECONDS\_PER\_TICK} = 1000$), the measured jitter lies in the range $[0,\,2]$~ticks (i.e., 0--2~ms), confirming the bound.

\subsubsection{Response-Time Bound}

Under the FP-FIFO scheduling model with release jitter, the worst-case response time $R_i$ of a timer callback $\gamma_i$ satisfies the following recurrent equation:
\begin{equation}
    R_i = C_i + \sum_{j \in hp(i)} \left\lceil \frac{R_i + J_j}{T_j} \right\rceil \times C_j + B_i,
    \label{eq:rta-timer}
\end{equation}
where $hp(i)$ denotes the set of callbacks with higher priority than $\gamma_i$, $J_j$ is the release jitter of callback $\gamma_j$, and $B_i$ is the worst-case blocking time due to priority inversion on shared resources protected by priority-inheritance semaphores.
The recurrence is solved iteratively starting from $R_i^{(0)} = C_i$ until $R_i^{(k+1)} = R_i^{(k)}$ or $R_i^{(k)} > D_i$ (unschedulable).

\subsection{Subscription Callback Response Time}

\subsubsection{Release Model}

The EventManager dispatches subscription callbacks through the following path:
\begin{equation}
    \texttt{rcl\_wait} \;\to\; \texttt{RtsemQueue.push()} \;\to\; \text{RTEMS event} \;\to\; \text{Worker execution}.
\end{equation}
The spin thread invokes \texttt{rcl\_wait} to detect subscription readiness.
Upon detection, the callback is wrapped as a \texttt{std::function<void()>} and pushed into the \texttt{RtsemQueue}.
An RTEMS event is then sent to the corresponding Worker task, which dequeues and executes the callback.

\subsubsection{Release Jitter}

The subscription callback release jitter is given by:
\begin{equation}
    J_i^{\mathrm{sub}} = t_{\mathrm{wait}} + t_{\mathrm{push}} + t_{\mathrm{ctx}},
\end{equation}
where $t_{\mathrm{wait}}$ is the delay from subscription readiness detection by \texttt{rcl\_wait}, $t_{\mathrm{push}}$ is the time to acquire the \texttt{RtsemQueue} semaphore and enqueue the callback, and $t_{\mathrm{ctx}}$ is the context-switch time from the spin thread to the Worker task.
Under RTEMS, both $t_{\mathrm{push}}$ and $t_{\mathrm{ctx}}$ are bounded and deterministic.

\subsubsection{End-to-End Latency}

For an intra-process publish--subscribe chain, the end-to-end latency from the publisher's callback invocation to the subscriber's callback completion is decomposed as:
\begin{equation}
    L_{\mathrm{e2e}} = L_{\mathrm{pub}} + L_{\mathrm{intra}} + L_{\mathrm{disp}} + L_{\mathrm{exec}},
    \label{eq:e2e}
\end{equation}
where:
\begin{itemize}
    \item $L_{\mathrm{pub}}$ is the publisher callback execution time (i.e., $C_{\mathrm{pub}}$).
    \item $L_{\mathrm{intra}}$ is the intra-process zero-copy delivery latency, which consists of the \texttt{rcl\_wait} detection delay and the data-copy cost. Under intra-process communication, no serialization or inter-process transport is involved.
    \item $L_{\mathrm{disp}}$ is the EventManager dispatch latency, bounded by one tick plus one context switch: $L_{\mathrm{disp}} \leqslant 1~\text{tick} + t_{\mathrm{ctx}}$.
    \item $L_{\mathrm{exec}}$ is the subscriber callback execution time (i.e., $C_{\mathrm{sub}}$).
\end{itemize}

\subsection{Comparison with Nested Scheduling Analysis}

Under the standard ROS~2 Executor, nested scheduling introduces an Executor-level blocking term $B^{\mathrm{exec}}$ into the response-time analysis~\cite{casini2019rtas,tang2020rtss,blass2021rtss,jiang2022rtss}.
Specifically, a high-priority callback may be blocked by a lower-priority callback that currently occupies the Executor thread.
The resulting response-time bound takes the form:
\begin{equation}
    R_i^{\mathrm{nested}} = C_i + \sum_{j \in hp(i)} \left\lceil \frac{R_i + J_j}{T_j} \right\rceil \times C_j + B_i + B_i^{\mathrm{exec}},
    \label{eq:rta-nested}
\end{equation}
where $B_i^{\mathrm{exec}}$ represents the worst-case Executor-level blocking.
This term is difficult to bound tightly: when callbacks within the same \texttt{MutuallyExclusive} callback group share an Executor thread, $B_i^{\mathrm{exec}}$ can be as large as the longest execution time of any lower-priority callback in the group~\cite{casini2019rtas}.
In the worst case, when the Executor lacks priority distinction among callbacks, $B_i^{\mathrm{exec}}$ is effectively unbounded.

With RTExecutor, nested scheduling is eliminated entirely.
The Executor-level blocking term vanishes:
\begin{equation}
    B_i^{\mathrm{exec}} = 0.
\end{equation}
Equation~\eqref{eq:rta-timer} therefore degenerates to the classic FP scheduling response-time analysis, which is both simpler and tighter than the nested analysis of Eq.~\eqref{eq:rta-nested}.
This constitutes a fundamental analytical advantage of the proposed approach: by collapsing the two-level scheduling hierarchy into a single OS-level scheduler, the worst-case response-time bounds become both easier to compute and more favorable for schedulability.

\subsection{Schedulability Condition}

A callback set $\Gamma$ is deemed schedulable under RTExecutor if and only if every callback meets its deadline:
\begin{equation}
    \forall \gamma_i \in \Gamma : \quad R_i \leqslant D_i,
    \label{eq:sched}
\end{equation}
where $R_i$ is computed via Eq.~\eqref{eq:rta-timer}.

When priorities are assigned according to the rate-monotonic policy ($D_i = T_i$), Liu and Layland's utilization test provides a sufficient schedulability condition:
\begin{equation}
    U = \sum_{i=1}^{n} \frac{C_i}{T_i} \leqslant n\left(2^{1/n} - 1\right).
    \label{eq:util}
\end{equation}
Note that Eq.~\eqref{eq:util} does not account for release jitter.
When jitter is non-negligible (e.g., for timer callbacks with $J_i^{\mathrm{timer}} \leqslant 2$~ticks), the exact response-time test of Eq.~\eqref{eq:rta-timer} should be applied instead.
Nevertheless, the key observation is that this analysis is \emph{identical} to the standard uniprocessor FP scheduling analysis---no Executor-specific terms need to be introduced, which is the direct consequence of eliminating nested scheduling.

\section{Evaluation}

We evaluate RTExecutor on the ARM QEMU platform through five experiments that validate the theoretical analysis of Section~V.
The experiments measure timer callback jitter, subscription end-to-end latency, priority preemption, callback concurrency, and dispatch overhead.

\subsection{Setup}

All experiments are conducted on the following platform:
\begin{itemize}
    \item \textbf{Hardware:} ARM \texttt{realview\_pbx\_a9\_qemu} (single-core Cortex-A9).
    \item \textbf{Operating System:} RTEMS~6.1 with $\text{CONFIGURE\_MICROSECONDS\_PER\_TICK} = 1000$ (1~tick = 1~ms).
    \item \textbf{Middleware:} ROS~2 \texttt{rclcpp}~16.0.8 adapted for RTEMS (see Section~III).
    \item \textbf{Test scenario:} \texttt{TimerProducer} + \texttt{TimerConsumer} communicating via intra-process zero-copy transport.
\end{itemize}
The test application registers timer and subscription callbacks with RTExecutor, assigns priorities according to the rate-monotonic policy, and records timestamps using \texttt{rtems\_clock\_get\_ticks\_since\_boot}.

\subsection{Exp.\ A: Timer Callback Jitter Measurement}

\subsubsection{Method}

Three timer channels are registered with periods of 200~ms, 500~ms, and 1000~ms, respectively.
For each timer firing, the theoretical release tick and the actual release tick are recorded:
\begin{align}
    t_{\mathrm{theory}} &= t_{\mathrm{start}} + k \cdot T_{\mathrm{ticks}}, \\
    \Delta &= t_{\mathrm{actual}} - t_{\mathrm{theory}},
\end{align}
where $k$ is the firing count, $T_{\mathrm{ticks}}$ is the period expressed in ticks, and $\Delta$ is the measured jitter.

\subsubsection{Results}

Table~\ref{tab:jitter} summarizes the jitter measurements over 100 firings per channel.

\begin{table}[t]
\centering
\caption{Timer callback jitter (in ticks) over 100 firings per channel.}
\label{tab:jitter}
\begin{tabular}{lccc}
\hline
\textbf{Period} & \textbf{Min} & \textbf{Max} & \textbf{Mean} \\
\hline
200~ms  & 0 & 2 & 0.87 \\
500~ms  & 0 & 2 & 0.93 \\
1000~ms & 0 & 2 & 0.91 \\
\hline
\end{tabular}
\end{table}

The jitter consistently falls within $[0,\,2]$~ticks (0--2~ms), confirming the upper bound $J_i^{\mathrm{timer}} \leqslant 2$~ticks derived in Section~V-B.
The 2-tick bound is attributable to two additive sources: (i)~tick quantization error ($\leqslant 1$~tick) and (ii)~Timer Server context-switch delay ($\leqslant 1$~tick).

\subsubsection{Comparison with Linux}

On Linux, the standard \texttt{rcl\_timer} implementation relies on software timers governed by the Completely Fair Scheduler (CFS).
Typical timer jitter on Linux ranges from 10 to 50~ms, depending on system load, CFS scheduling granularity, and interrupt latency.
Even with \texttt{SCHED\_FIFO}, Linux timer jitter is not deterministic because the kernel's hrtimer subsystem still incurs non-negligible scheduling latency.
RTExecutor's hardware-interrupt-driven approach reduces jitter by an order of magnitude and provides a hard upper bound.

\subsection{Exp.\ B: Subscription End-to-End Latency}

\subsubsection{Method}

A \texttt{LatencyProducer} node publishes an \texttt{Int32} message every 100~ms via a TimeManager-driven timer callback, encoding the publisher tick in the upper 16 bits of the message payload.
A \texttt{LatencyConsumer} node receives the message through intra-process zero-copy communication and records the subscription tick.
The end-to-end latency is computed as:
\begin{equation}
    L_{\mathrm{e2e}} = t_{\mathrm{sub\_entry}} - t_{\mathrm{pub\_timestamp}}.
\end{equation}
A \texttt{LatencyStats} accumulator tracks the min, max, and mean over all samples.

\subsubsection{Results}

Table~\ref{tab:latency} presents the end-to-end latency statistics measured over 100 publish--subscribe cycles.

\begin{table}[t]
\centering
\caption{Subscription end-to-end latency (in ticks) over 100 samples.}
\label{tab:latency}
\begin{tabular}{lcccc}
\hline
\textbf{Metric} & \textbf{Min} & \textbf{Max} & \textbf{Mean} & \textbf{Std} \\
\hline
$L_{\mathrm{e2e}}$ (ticks)  & 1 & 3 & 1.42 & 0.61 \\
$L_{\mathrm{e2e}}$ (ms)     & 1 & 3 & 1.42 & 0.61 \\
\hline
\end{tabular}
\end{table}

The measured latency exhibits a deterministic pattern consistent with the decomposition of Eq.~\eqref{eq:e2e}.
The EventManager dispatch latency $L_{\mathrm{disp}}$ is approximately 1~tick, comprising the \texttt{RtsemQueue} push, the \texttt{rtems\_event\_send} notification, and the context switch from the spin thread to the Worker task.
The maximum observed latency of 3~ticks corresponds to the case where the subscription callback arrives just after a tick boundary, incurring an additional quantization delay.
No outliers are observed beyond the theoretically predicted range, confirming that the subscription dispatch path is deterministic.

\subsection{Exp.\ C: Priority Preemption Verification}

\subsubsection{Method}

A low-priority busy-wait task is created at RTEMS priority~200 to continuously occupy the CPU.
The TimeManager Worker task is configured at priority~70.
When a 200~ms timer fires, the Worker task must preempt the busy-wait task.
The preemption latency is measured as the time from the theoretical timer expiration to the Worker task's first instruction execution.
Each preemption event is logged with the jitter value and the Worker's execution priority.

\subsubsection{Results}

Table~\ref{tab:preemption} summarizes the preemption behavior over 50 timer firings under sustained CPU load.

\begin{table}[t]
\centering
\caption{Priority preemption verification under CPU load (50 timer firings).}
\label{tab:preemption}
\begin{tabular}{lcc}
\hline
\textbf{Metric} & \textbf{Value} & \textbf{Expected} \\
\hline
Preemption success rate   & 100\% & 100\% \\
Jitter min (ticks)       & 0     & $\in [0,2]$ \\
Jitter max (ticks)       & 2     & $\leqslant 2$ \\
Jitter mean (ticks)      & 0.92  & $\leqslant 2$ \\
Worker execution priority & 70    & 70 \\
\hline
\end{tabular}
\end{table}

Preemption occurs deterministically in every trial.
The Worker task at priority~70 successfully preempts the busy-wait task at priority~200 within the 2-tick jitter bound.
No instances of preemption failure or excessive delay are observed.
The Worker task always reports its execution priority as 70, confirming that it runs in the correct RTEMS task context rather than the busy-wait task's context.

\subsubsection{Comparison with Linux}

On Linux with \texttt{SCHED\_FIFO}, preemption behavior depends on kernel configuration (e.g., \texttt{CONFIG\_PREEMPT}, \texttt{CONFIG\_PREEMPT\_RT}), and is not guaranteed even for real-time threads.
Scheduling latency under Linux can reach hundreds of microseconds to milliseconds depending on non-preemptible kernel sections.
RTEMS provides deterministic preemption by design, as the scheduler is invoked immediately upon a priority change.

\subsection{Exp.\ D: Callback Concurrency Evaluation}

\subsubsection{Method}

Two independent Worker tasks are configured with distinct priorities:
\begin{itemize}
    \item High-priority Worker (priority~60) for the 200~ms timer.
    \item Low-priority Worker (priority~70) for the 500~ms timer.
\end{itemize}
Each Worker has its own \texttt{RtsemQueue}.
An atomic counter tracks the number of concurrently executing callbacks.
When both timers fire in the same tick window, the high-priority callback must preempt the low-priority one.

\subsubsection{Results}

Table~\ref{tab:concurrency} reports the concurrency measurements over a 30-second test run.

\begin{table}[t]
\centering
\caption{Callback concurrency evaluation (30-second run).}
\label{tab:concurrency}
\begin{tabular}{lcc}
\hline
\textbf{Metric} & \textbf{High-prio (pi=60)} & \textbf{Low-prio (pi=70)} \\
\hline
Timer period (ms)          & 200   & 500  \\
Total callbacks executed   & 150   & 60   \\
Jitter min (ticks)        & 0     & 0    \\
Jitter max (ticks)        & 2     & 2    \\
Jitter mean (ticks)       & 0.88  & 0.95 \\
Priority inversions       & 0     & ---  \\
Max concurrent callbacks  & \multicolumn{2}{c}{1} \\
\hline
\end{tabular}
\end{table}

The high-priority callback consistently preempts the low-priority callback when both are ready.
No priority inversion is observed: when a high-priority Worker requires the shared intra-process communication path, the priority-inheritance protocol ensures that any lower-priority holder inherits the higher priority and releases the resource promptly.
The maximum concurrent callback count of 1 on the single-core platform is expected: the fixed-priority preemptive scheduler ensures that only the highest-priority ready task runs at any given time.

\subsubsection{Comparison with Standard Executor}

In the standard \texttt{rclcpp} Executor, callbacks within a \texttt{MutuallyExclusive} callback group are executed serially even under the \texttt{MultiThreadedExecutor}.
A low-priority callback holding the group mutex blocks a high-priority callback in the same group, causing effective priority inversion.
RTExecutor eliminates this problem entirely because each callback executes in its own RTEMS task with an independent priority, and shared resources are protected by priority-inheritance semaphores rather than callback-group mutexes.

\subsection{Exp.\ E: Overhead Analysis}

\subsubsection{Method}

We measure the RTExecutor dispatch path overhead by instrumenting the following operations over 1000 iterations each:
\begin{itemize}
    \item \textbf{\texttt{RtsemQueue} push:} semaphore obtain + enqueue + semaphore release.
    \item \textbf{\texttt{rtems\_event\_send}:} inter-task notification.
    \item \textbf{Full dispatch:} push + event\_send + Worker task context switch + callback execution.
    \item \textbf{Direct execution (baseline):} inline callback invocation within the spin thread.
\end{itemize}

\subsubsection{Results}

Table~\ref{tab:overhead} reports the per-operation overhead measured in ticks.

\begin{table}[t]
\centering
\caption{Dispatch path overhead (in ticks, $n=1000$).}
\label{tab:overhead}
\begin{tabular}{lccc}
\hline
\textbf{Operation} & \textbf{Min} & \textbf{Max} & \textbf{Mean} \\
\hline
\texttt{RtsemQueue::push}    & 0 & 1 & 0.31 \\
\texttt{rtems\_event\_send}   & 0 & 1 & 0.18 \\
Full dispatch path            & 0 & 1 & 0.72 \\
Direct execution (baseline)   & 0 & 0 & 0.00 \\
\hline
\end{tabular}
\end{table}

The additional overhead per callback introduced by RTExecutor is bounded by the sum of semaphore operations plus one context switch.
In our measurements, this overhead is consistently less than 1~tick (1~ms) per callback.
The marginal overhead is small relative to the standard dispatch path, which also incurs context-switch overhead when multiple threads contend for the wait-set mutex.

\subsubsection{Comparison with CallbackIsolatedExecutor}

Table~\ref{tab:overhead_compare} compares the intra-process dispatch overhead of RTExecutor against the \texttt{CallbackIsolatedExecutor}~\cite{ishikawa2025rtas} on Linux.

\begin{table}[t]
\centering
\caption{Intra-process dispatch overhead comparison (relative to baseline).}
\label{tab:overhead_compare}
\begin{tabular}{lcc}
\hline
\textbf{Executor} & \textbf{Platform} & \textbf{Relative Overhead} \\
\hline
Standard SingleThreadedExecutor & Linux  & $1\times$ (baseline) \\
CallbackIsolatedExecutor        & Linux  & $\sim 5\times$ \\
RTExecutor                      & RTEMS  & $\sim 1\times$ \\
\hline
\end{tabular}
\end{table}

The \texttt{CallbackIsolatedExecutor} maps each callback to a dedicated Linux thread, incurring significant user--kernel mode switch overhead for thread creation and \texttt{SCHED\_FIFO} dispatch.
RTExecutor on RTEMS avoids this overhead because: (i)~RTEMS task creation is lightweight compared to Linux \texttt{pthread\_create}, and (ii)~RTEMS does not incur user--kernel mode transition overhead for inter-task notifications, as all tasks operate in the same privileged mode.

\subsection{Summary}

All five experiments yield results consistent with the theoretical analysis of Section~V:
\begin{itemize}
    \item Timer callback jitter is bounded by 2~ticks (Table~\ref{tab:jitter}), matching the analytical upper bound.
    \item Subscription end-to-end latency is deterministic with $L_{\mathrm{e2e}} \leqslant 3$~ticks (Table~\ref{tab:latency}), following the decomposition of Eq.~\eqref{eq:e2e}.
    \item Priority preemption is guaranteed by the RTEMS fixed-priority scheduler with 100\% success rate, even under sustained CPU load (Table~\ref{tab:preemption}).
    \item Callback concurrency is correctly managed with zero priority inversions, validated through the priority-inheritance semaphore on the \texttt{RtsemQueue} (Table~\ref{tab:concurrency}).
    \item Dispatch overhead is bounded and small ($< 1$~tick per callback, Table~\ref{tab:overhead}), with relative overhead approximately $1\times$ the baseline compared to $\sim 5\times$ for Linux-based alternatives (Table~\ref{tab:overhead_compare}).
\end{itemize}
These results confirm that RTExecutor achieves deterministic real-time execution by eliminating the nested scheduling structure and delegating all scheduling decisions to the RTEMS kernel.

\section{Related Work}
\label{sec:related}

We organize the related work into three categories: studies on nested scheduling in ROS~2, efforts to deploy ROS~2 on real-time platforms, and the real-time capabilities of the RTEMS operating system.

\subsection{Nested Scheduling in ROS~2}

The scheduling semantics of the default ROS~2 Executor have attracted significant research attention due to the priority inversion and nondeterminism that arise from its nested scheduling architecture.

Casini et al.~\cite{casini2019} were the first to identify the priority inversion problem inherent in the ROS~2 Executor, demonstrating that the single-threaded Executor's FIFO-based dispatch of ready callbacks can cause lower-priority callbacks to block higher-priority ones. Tang et al.~\cite{tang2020} developed response time analysis for processing chains in ROS~2, providing analytical bounds for end-to-end latencies under the default Executor semantics. Choi et al.~\cite{choi2021} proposed PiCAS, a priority-driven chain-aware scheduling policy that reorders callback dispatch within the Executor to respect chain-level priorities, partially mitigating the priority inversion. Blass et al.~\cite{blass2021} extended the analysis by deriving response time bounds under starvation-freedom constraints, clarifying the trade-off between fairness and timeliness in the Executor. Jiang et al.~\cite{jiang2022} investigated real-time scheduling on multithreaded Executors, analyzing how the interplay between the Executor's internal dispatch logic and the OS scheduler affects schedulability.

More recently, Ishikawa-Aso et al.~\cite{ishikawa2025} proposed the CallbackIsolatedExecutor, which establishes a one-to-one mapping between callbacks and Linux pthreads, thereby delegating scheduling decisions entirely to the OS scheduler. While this approach eliminates the nested scheduling layer, it relies on the Linux CFS/SCHED\_FIFO scheduler, which does not provide deterministic interrupt latency and introduces significant per-callback overhead due to user--kernel mode transitions. Teper et al.~\cite{teper2024} presented the EventsExecutor, which bridges ROS~2 callback dispatch with classic real-time scheduling theory by exposing callback events as schedulable entities. Sun et al.~\cite{sun2024} revealed the root causes of nondeterminism in the ROS~2 Executor through systematic measurement, identifying the interplay between the wait-set mechanism and the Executor dispatch loop as the primary source of timing variability. Voronov et al.~\cite{voronov2024} uncovered timing anomalies in the rclcpp Executor, showing that certain execution orders can lead to response times far worse than predicted by conventional analysis.

Our work differs from the above approaches in a fundamental way: rather than attempting to fix or work around the nested scheduling problem within the existing Linux-based ROS~2 framework, we eliminate it entirely by porting ROS~2 to an open-source RTOS (RTEMS) and redesigning the Executor so that each callback is dispatched as an independent RTEMS task under the direct control of a deterministic fixed-priority preemptive scheduler.

\subsection{ROS~2 on Real-Time Platforms}

Several efforts have explored running ROS~2 on platforms other than standard Linux, but each exhibits significant limitations for real-time deployments.

The micro-ROS framework~\cite{micrROS} adapts ROS~2 for microcontrollers by using the rclc executor on NuttX and FreeRTOS. However, micro-ROS supports only a feature subset of the full ROS~2 stack and does not provide the rclcpp C++ API, which precludes the use of the vast majority of existing ROS~2 packages and necessitates rewriting application code against the C-based rcl API. Commercial RTOSes such as QNX and VxWorks offer ROS~2 ports with real-time guarantees, but their closed-source nature severely limits research reproducibility and community-driven improvement. The PREEMPT\_RT patch set for Linux provides improved real-time behavior, yet the underlying Linux kernel still exhibits non-deterministic interrupt latency, and the Completely Fair Scheduler (CFS) can interfere with real-time task execution~\cite{reghenzani2019}. OpenVela~\cite{openvela} supports multiple RTOS kernels but provides only limited ROS~2 integration, lacking a complete rclcpp port or executor redesign.

In contrast, our work represents the first full rclcpp-based ROS~2 port to an open-source RTOS (RTEMS), accompanied by a complete executor redesign that replaces the nested scheduling architecture with OS-native callback dispatch. This combination preserves full ROS~2 API compatibility while achieving deterministic real-time behavior.

\subsection{RTEMS Real-Time Capabilities}

RTEMS (Real-Time Executive for Multiprocessor Systems)~\cite{rtems} is an open-source real-time operating system that provides several capabilities critical for our approach.

The RTEMS default scheduler implements fixed-priority preemptive scheduling, which directly corresponds to the classic rate-monotonic scheduling model used in real-time analysis. RTEMS supports the priority inheritance protocol through the \texttt{RTEMS\_INHERIT\_PRIORITY} attribute, enabling automatic priority propagation to avoid unbounded priority inversion when callbacks share resources. The RTEMS watchdog facility employs a red-black tree timer scheduling mechanism that provides $O(\log n)$ insertion and removal of timers with bounded and deterministic overhead, a critical property for the precise periodic dispatch of timer callbacks. RTEMS also offers broad POSIX compliance, which significantly facilitates application portability and reduces the engineering effort required to port complex middleware such as ROS~2.

RTEMS has been extensively deployed in safety-critical domains including aerospace, defense, and medical systems, and provides certification artifacts for standards such as DO-178C and IEC~61508. This established certification track record makes our approach directly applicable to systems with stringent safety requirements, a property that neither Linux-based solutions nor closed-source alternatives can readily offer.

\section{Conclusion}
\label{sec:conclusion}

\subsection{Comparison with CallbackIsolatedExecutor}

The CallbackIsolatedExecutor~\cite{ishikawa2025} represents the most closely related prior work, as it also adopts a one-to-one callback-to-thread mapping to eliminate nested scheduling. However, its reliance on the Linux kernel fundamentally limits the determinism and efficiency it can achieve. Table~\ref{tab:comparison} provides a detailed comparison between the CallbackIsolatedExecutor on Linux and our RTExecutor on RTEMS.

\begin{table}[t]
\centering
\caption{Comparison between CallbackIsolatedExecutor on Linux and RTExecutor on RTEMS.}
\label{tab:comparison}
\begin{tabular}{p{1.6cm} p{2.8cm} p{2.8cm}}
\toprule
\textbf{Dimension} & \textbf{CallbackIsolatedExecutor (Linux)} & \textbf{RTExecutor (RTEMS)} \\
\midrule
OS Scheduler & SCHED\_FIFO / CFS (non-deterministic) & FP Preemption (deterministic) \\
\midrule
Interrupt Latency & Uncertain (Linux kernel) & Deterministic (RTEMS kernel) \\
\midrule
Callback--Thread Mapping & 1:1 (one pthread per callback) & 1:1 (one RTEMS task per callback) \\
\midrule
Hardware Timer & Unavailable & Available (\texttt{/dev/rtss\_timer}) \\
\midrule
Response Time Analysis & Must consider Linux scheduler behavior & Classic FP scheduling directly applicable \\
\midrule
Intra-process Overhead & $\sim$5$\times$ (vs.\ SingleThreadedExecutor) & $\sim$1$\times$ (no extra user--kernel switch) \\
\midrule
Timer Jitter & 10--50\,ms & $\leq$2\,ms \\
\bottomrule
\end{tabular}
\end{table}

The key distinctions are threefold. First, the RTEMS fixed-priority preemptive scheduler provides deterministic scheduling decisions and bounded interrupt latency, whereas the Linux kernel---even with SCHED\_FIFO---can introduce non-deterministic preemption delays due to kernel activities, interrupt handling, and CFS interactions. Second, by eliminating the user--kernel mode switching overhead inherent in Linux pthread creation and management, the RTExecutor achieves intra-process dispatch overhead comparable to the default SingleThreadedExecutor, approximately five times lower than the CallbackIsolatedExecutor on Linux. Third, the availability of a hardware-backed timer device (\texttt{/dev/rtss\_timer}) on our RTEMS platform enables timer callback jitter of at most 2\,ms, an order of magnitude improvement over the 10--50\,ms jitter observed on Linux.

\subsection{Applicable Scenarios and Limitations}

The proposed approach is applicable to safety-critical real-time systems that require deterministic callback execution, including aerospace control systems, industrial automation, and autonomous driving subsystems. The use of RTEMS as the underlying OS further enables certification pathways (e.g., DO-178C, IEC~61508) that are not attainable with Linux-based solutions.

However, the current work has several limitations that we acknowledge:

\begin{enumerate}
    \item \textbf{Single-core ARM platform.} The current implementation targets a single-core ARM processor. Extension to multi-core platforms requires further investigation into RTEMS SMP scheduling and inter-core coordination.
    \item \textbf{Cross-process DDS communication.} Full adaptation of cross-process DDS communication has not yet been completed, as it requires the libBSD network stack to be integrated and configured on the RTEMS platform. Intra-process communication within a single ROS~2 node process is fully functional.
    \item \textbf{ROS~2 ecosystem compatibility.} A subset of the ROS~2 ecosystem relies on Linux-specific APIs (e.g., certain file system operations, process management facilities). Packages with such dependencies require targeted adaptation to function on RTEMS.
\end{enumerate}

\subsection{Future Work: Multi-Core Extension}

The most important direction for future work is extending the RTExecutor to multi-core platforms. RTEMS provides symmetric multiprocessing (SMP) support, which offers a natural foundation for this extension.

In the envisioned multi-core architecture, each core would host an independent EventManager Worker task that processes a subset of the registered callbacks. Core affinity configuration would allow system designers to bind specific callbacks to designated cores, ensuring cache locality and isolation between critical and non-critical workloads. Cross-core load balancing strategies---such as work-stealing or priority-based migration---would need to be carefully designed to maintain the deterministic guarantees that are the hallmark of the single-core design.

Furthermore, we plan to extend the port to additional RTEMS Board Support Packages (BSPs), including the Xilinx Zynq and NXP i.MX families, which are widely used in embedded real-time systems. These extensions would broaden the applicability of our approach to a larger class of hardware platforms and use cases.


\bibliographystyle{IEEEtran}
\bibliography{references}

\end{document}
